% ==============================================================================
% PLANTILLA_ADULTO.TEX — MACRO Y DEFINICIÓN DE FICHA PARA TUTORES Y FAMILIAS
% Asignatura: Metodologías de Desarrollo Ágil (MDA) — Universidad de Granada
% Proyecto AUTOVIDA (26-175) | Reto CeRA — Agenda Accesible PTVAL
% ==============================================================================
% Campos canónicos: 11 (Con Occupation y Education, sin Company Information)
% Uso: % ==============================================================================
% PLANTILLA_ADULTO.TEX — MACRO Y DEFINICIÓN DE FICHA PARA TUTORES Y FAMILIAS
% Asignatura: Metodologías de Desarrollo Ágil (MDA) — Universidad de Granada
% Proyecto AUTOVIDA (26-175) | Reto CeRA — Agenda Accesible PTVAL
% ==============================================================================
% Campos canónicos: 11 (Con Occupation y Education, sin Company Information)
% Uso: % ==============================================================================
% PLANTILLA_ADULTO.TEX — MACRO Y DEFINICIÓN DE FICHA PARA TUTORES Y FAMILIAS
% Asignatura: Metodologías de Desarrollo Ágil (MDA) — Universidad de Granada
% Proyecto AUTOVIDA (26-175) | Reto CeRA — Agenda Accesible PTVAL
% ==============================================================================
% Campos canónicos: 11 (Con Occupation y Education, sin Company Information)
% Uso: % ==============================================================================
% PLANTILLA_ADULTO.TEX — MACRO Y DEFINICIÓN DE FICHA PARA TUTORES Y FAMILIAS
% Asignatura: Metodologías de Desarrollo Ágil (MDA) — Universidad de Granada
% Proyecto AUTOVIDA (26-175) | Reto CeRA — Agenda Accesible PTVAL
% ==============================================================================
% Campos canónicos: 11 (Con Occupation y Education, sin Company Information)
% Uso: \input{plantilla_adulto.tex} y luego invocar \fichaAdulto[...] o \fichaAdultoBlanco
% ==============================================================================

\ifdefined\isSubFile
    % Si se incluye desde otro documento, los paquetes ya están cargados
\else
    \documentclass[10pt]{article}
    \usepackage[utf8]{inputenc}
    \usepackage[spanish,es-tabla]{babel}
    \usepackage[margin=0.8cm,top=0.6cm,bottom=0.6cm]{geometry}
    \usepackage{helvet}
    \renewcommand{\familydefault}{\sfdefault}
    \usepackage{microtype}
    \usepackage{graphicx}
    \usepackage{xcolor}
    \usepackage{tikz}
    \usepackage{enumitem}
    \usepackage{fontawesome5}
    \usepackage{tcolorbox}
    \tcbuselibrary{skins}
    \pagestyle{empty}
\fi

% Colores del sistema
\providecolor{primaryColor}{RGB}{20, 65, 105}     % Azul corporativo
\providecolor{secondaryColor}{RGB}{35, 115, 170}  % Azul medio accesible
\providecolor{tealColor}{RGB}{20, 130, 140}       % Verde azulado
\providecolor{greenColor}{RGB}{30, 135, 75}       % Verde bosque
\providecolor{accentColor}{RGB}{215, 105, 25}     % Naranja cálido
\providecolor{purpleColor}{RGB}{120, 60, 150}     % Púrpura suave
\providecolor{dangerColor}{RGB}{185, 45, 35}      % Rojo criticidad
\providecolor{bgCard}{RGB}{252, 253, 255}         % Fondo limpio
\providecolor{borderCard}{RGB}{215, 225, 235}     % Borde sutil
\providecolor{textDark}{RGB}{30, 40, 50}          % Texto principal oscuro
\providecolor{textMuted}{RGB}{95, 110, 125}       % Texto secundario

% Configuración de listas
\setlist[itemize]{leftmargin=3.0mm, itemsep=0.4pt, parsep=0pt, topsep=0.4pt, partopsep=0pt}

% Estilos de caja fija para Adulto
\tcbset{
    cardAdultoBox/.style={
        enhanced,
        boxrule=0.65pt,
        arc=1.8mm,
        height=27.2mm,
        left=5.0pt, right=5.0pt, top=1.5pt, bottom=1.5pt,
        boxsep=0pt,
        colback=bgCard,
        colframe=borderCard,
        fonttitle=\bfseries\small,
        fontupper=\fontsize{8.2pt}{9.6pt}\selectfont,
    },
    cardAdultoWideBox/.style={
        enhanced,
        boxrule=0.65pt,
        arc=1.8mm,
        height=20.0mm,
        left=6.0pt, right=6.0pt, top=2.0pt, bottom=2.0pt,
        boxsep=0pt,
        colback=bgCard,
        colframe=borderCard,
        fonttitle=\bfseries\small,
    }
}

% Macro para badges visuales (si no existe)
\providecommand{\badge}[2]{%
    \tikz[baseline=(b.base)] \node[fill=#1!12, text=#1!90!black, draw=#1!35, line width=0.4pt, rounded corners=1.2mm, inner sep=2.0pt, font=\bfseries\scriptsize] (b) {#2};%
}

% Definición de claves con pgfkeys para \fichaAdulto
\pgfkeys{
    /fichaAdulto/.is family, /fichaAdulto,
    default/.style = {
        bannerTitulo = {\faChalkboardTeacher\ FICHA DE PERSONA: TUTORES Y FAMILIAS (METODOLOGÍA DCU)},
        subtitulo = {Proyecto AUTOVIDA (26-175) \textbar\ Reto CeRA — Agenda Accesible PTVAL \textbar\ Asignatura MDA (UGR)},
        equipo = {\textbf{Equipo CeRA-MDA} (5 integrantes)\\\ Course 2026/2027 \textbar\ \textbf{Práctica 1}},
        nombre = {[Nombre y Apellidos del Usuario]},
        nombreCorto = {[ADULTO]},
        edad = {[Edad]},
        icono = {\faChalkboardTeacher},
        foto = {},
        badgeRol = {[Rol: Tutor / Familiar / Admin]},
        badgeTarea = {[Ámbito de Actuación]},
        centro = {CEE Purísima Concepción},
        cita = {[Cita textual representativa que exprese sus objetivos, inquietudes de gestión o necesidades respecto a la app CeRA].},
        demographic = {
            \item \textbf{Edad y Género:} [[Edad] años, [Género]].
            \item \textbf{Ubicación:} [[Ciudad / Zona residencial]].
        },
        occupation = {
            \item \textbf{Puesto actual:} [[Puesto docente, profesión o rol familiar]].
            \item \textbf{Antigüedad:} [[Años en el centro/entidad] \textbar\ [Años en puesto]].
            \item \textbf{Responsabilidades:} [[Supervisión, docencia, apoyo en el hogar]].
        },
        education = {
            \item \textbf{Titulación (Degree):} [[Nivel formativo o titulación alcanzada]].
            \item \textbf{Especialidad (Major):} [[Educación Especial, Pedagogía, etc.]].
            \item \textbf{Cursos realizados:} [[Formación en NEAE, SAAC, TIC o gestión]].
        },
        computerExp = {
            \item \textbf{Habilidades informáticas:} [[Nivel de competencia digital: medio/avanzado]].
            \item \textbf{Años de experiencia:} [[Años de uso regular de dispositivos]].
            \item \textbf{Hábitos de uso:} [[Navegación web, apps de gestión, smartphones]].
        },
        productExp = {
            \item \textbf{Productos similares:} [[Otras herramientas usadas, papel, agendas]].
            \item \textbf{Limitaciones actuales:} [[Sobrecarga de trabajo, falta de personalización]].
            \item \textbf{Expectativas:} [[Ahorro de tiempo y canal de dudas]].
        },
        tasks = {
            \item \textbf{Tareas Primarias:} 1) [[Configurar tareas]] $\rightarrow$ 2) [[Asignar a alumnos]] $\rightarrow$ 3) [[Supervisar ejecución]] $\rightarrow$ 4) [[Gestionar comedor/avisos]].
            \item \textbf{Tareas Secundarias:} [[Revisar estadísticas / chat con familias]].
        },
        domainKnowledge = {
            \item \textbf{Comprensión del área:} [[Conocimiento sobre necesidades NEAE y rutinas]].
            \item \textbf{Conocimiento del usuario:} [[Grado de personalización que requiere cada estudiante]].
        },
        technology = {
            \item \textbf{Hardware:} [[Ordenador portátil / tablet / móvil Android]].
            \item \textbf{Software:} [[Navegador web / app móvil]].
            \item \textbf{Herramientas:} [[Conexión a internet]].
        },
        attitudes = {
            \item \textbf{Preferencias:} [[Sencillez, rapidez operativa y fiabilidad en tiempo real]].
            \item \textbf{Miedo a la tecnología:} [[Preocupación por pérdida de datos o fallos de red]].
            \item \textbf{Valores:} [[Foco en el progreso y bienestar del alumnado]].
        },
        learningStyle = {
            \item \textbf{Perfil de aprendizaje:} [[\textbf{Práctico} / \textbf{Estructurado} / \textbf{Visual}]].
            \item \textbf{Canales óptimos:} [[Tutoriales cortos, interfaz intuitiva sin manual]].
            \item \textbf{Curva de adopción:} [[Inmediata si soluciona problemas reales]].
        },
        criticality = {
            \item \textbf{Consecuencias de un error:} \textit{Logístico:} [[¿Dietas erróneas, tareas desincronizadas?]] \textit{Pedagógico:} [[¿Confusión en el horario del alumno, estrés?]]
            \item \textbf{Tolerancia del software:} [[Alertas claras ante cambios sensibles, confirmación de guardado]].
        },
        escenario = {[[Describir brevemente en 3-4 líneas cómo accede al panel de la app, planifica o revisa las tareas de sus estudiantes y cómo la herramienta le facilita la jornada eliminando fricciones organizativas.]]},
        badgesAjustes = {
            \badge{tealColor}{\faDesktop\ [[Panel Web]]} \hspace{2mm}
            \badge{primaryColor}{\faCalendarCheck\ [[Planificación]]} \hspace{2mm}
            \badge{accentColor}{\faUtensils\ [[Comedor/Dietas]]} \hspace{2mm}
            \badge{greenColor}{\faCheckCircle\ [[Supervisión]]}%
        }
    },
    bannerTitulo/.store in = \fAdBannerTitulo,
    subtitulo/.store in = \fAdSubtitulo,
    equipo/.store in = \fAdEquipo,
    nombre/.store in = \fAdNombre,
    nombreCorto/.store in = \fAdNombreCorto,
    edad/.store in = \fAdEdad,
    icono/.store in = \fAdIcono,
    foto/.store in = \fAdFoto,
    badgeRol/.store in = \fAdBadgeRol,
    badgeTarea/.store in = \fAdBadgeTarea,
    centro/.store in = \fAdCentro,
    cita/.store in = \fAdCita,
    demographic/.store in = \fAdDemographic,
    occupation/.store in = \fAdOccupation,
    education/.store in = \fAdEducation,
    computerExp/.store in = \fAdComputerExp,
    productExp/.store in = \fAdProductExp,
    tasks/.store in = \fAdTasks,
    domainKnowledge/.store in = \fAdDomainKnowledge,
    technology/.store in = \fAdTechnology,
    attitudes/.store in = \fAdAttitudes,
    learningStyle/.store in = \fAdLearningStyle,
    criticality/.store in = \fAdCriticality,
    escenario/.store in = \fAdEscenario,
    badgesAjustes/.store in = \fAdBadgesAjustes,
}

% COMANDO PRINCIPAL \fichaAdulto
\newcommand{\fichaAdulto}[1][]{%
    \begingroup
    \pgfkeys{/fichaAdulto, default, #1}%
    % 1. BANNER SUPERIOR
    \begin{tcolorbox}[
        enhanced,
        colback=tealColor!90!black,
        colframe=tealColor!90!black,
        arc=1.6mm,
        left=7pt, right=7pt, top=2.5pt, bottom=2.5pt,
        boxsep=0pt, coltext=white
    ]
        \begin{minipage}[c]{0.70\textwidth}
            {\normalsize\textbf{\fAdBannerTitulo}}\\[0.2mm]
            {\tiny \fAdSubtitulo}
        \end{minipage}\hfill
        \begin{minipage}[c]{0.29\textwidth}
            \raggedleft\tiny
            \fAdEquipo
        \end{minipage}
    \end{tcolorbox}
    \vspace{0.7mm}
    % 2. TARJETA DE IDENTIDAD DEL ADULTO
    \begin{tcolorbox}[
        enhanced,
        colback=white,
        colframe=tealColor!70,
        arc=2.0mm,
        boxrule=0.8pt,
        left=6pt, right=6pt, top=2.8pt, bottom=2.8pt,
        boxsep=0pt,
        drop shadow=black!5
    ]
        \begin{minipage}[c]{0.10\textwidth}
            \centering
            \begin{tikzpicture}
                \ifx\fAdFoto\empty
                    \draw[fill=tealColor!8, draw=tealColor!30, line width=0.5pt, rounded corners=2mm] (0,0) rectangle (1.4, 1.55);
                    \node at (0.7, 0.95) {\color{tealColor}\scalebox{1.5}{\fAdIcono}};
                    \node[align=center] at (0.7, 0.28) {\tiny\textbf{\color{tealColor}\fAdNombreCorto}\\\fontsize{4pt}{5pt}\selectfont\color{textMuted}\fAdEdad};
                \else
                    \begin{scope}
                        \clip[rounded corners=2mm] (0,0) rectangle (1.45, 1.55);
                        \node[inner sep=0pt] at (0.725, 0.775) {\includegraphics[width=1.45cm,height=1.55cm,keepaspectratio=false]{\fAdFoto}};
                    \end{scope}
                    \draw[draw=tealColor!60, line width=0.8pt, rounded corners=2mm] (0,0) rectangle (1.45, 1.55);
                \fi
            \end{tikzpicture}
        \end{minipage}\hfill
        \begin{minipage}[c]{0.88\textwidth}
            \begin{minipage}[b]{0.68\textwidth}\raggedright
                {\large\textbf{\color{tealColor!90!black}\fAdNombre}} \hspace{2mm}
                \badge{tealColor}{\fAdBadgeRol} \badge{secondaryColor}{\fAdBadgeTarea}
            \end{minipage}\hfill
            \begin{minipage}[b]{0.31\textwidth}
                \raggedleft\small\textbf{\color{tealColor!90!black}\faSchool\ \fAdCentro}
            \end{minipage}\\[0.6mm]
            \tikz[baseline] \node[fill=tealColor!5, draw=tealColor!25, rounded corners=1.2mm, inner sep=2.5pt, text width=0.98\linewidth] {
                \small\itshape\color{tealColor!90!black}
                \faQuoteLeft\ ``\fAdCita''\ \faQuoteRight
            };
        \end{minipage}
    \end{tcolorbox}
    \vspace{0.7mm}
    % 3. LAS 11 CAJAS DE ADULTOS (5 FILAS X 2 COLS + 1 FILA COMPLETA)
    \noindent\begin{minipage}[t]{0.485\textwidth}
        \begin{tcolorbox}[cardAdultoBox, colframe=primaryColor!50, coltitle=primaryColor, title={\faIdCard\ 1. Demographic characteristics}]
            \fontsize{8.0pt}{9.4pt}\selectfont
            \begin{itemize}
                \fAdDemographic
            \end{itemize}
        \end{tcolorbox}
    \end{minipage}\hfill
    \begin{minipage}[t]{0.485\textwidth}
        \begin{tcolorbox}[cardAdultoBox, colframe=primaryColor!50, coltitle=primaryColor, title={\faUserTie\ 2. Occupation experience}]
            \fontsize{8.0pt}{9.4pt}\selectfont
            \begin{itemize}
                \fAdOccupation
            \end{itemize}
        \end{tcolorbox}
    \end{minipage}
    \vspace{0.7mm}
    \noindent\begin{minipage}[t]{0.485\textwidth}
        \begin{tcolorbox}[cardAdultoBox, colframe=tealColor!60, coltitle=tealColor!90!black, title={\faGraduationCap\ 3. Education}]
            \fontsize{8.0pt}{9.4pt}\selectfont
            \begin{itemize}
                \fAdEducation
            \end{itemize}
        \end{tcolorbox}
    \end{minipage}\hfill
    \begin{minipage}[t]{0.485\textwidth}
        \begin{tcolorbox}[cardAdultoBox, colframe=secondaryColor!60, coltitle=secondaryColor!90!black, title={\faLaptop\ 4. Computer experience}]
            \fontsize{8.0pt}{9.4pt}\selectfont
            \begin{itemize}
                \fAdComputerExp
            \end{itemize}
        \end{tcolorbox}
    \end{minipage}
    \vspace{0.7mm}
    \noindent\begin{minipage}[t]{0.485\textwidth}
        \begin{tcolorbox}[cardAdultoBox, colframe=secondaryColor!60, coltitle=secondaryColor!90!black, title={\faSearchPlus\ 5. Specific product experience}]
            \fontsize{8.0pt}{9.4pt}\selectfont
            \begin{itemize}
                \fAdProductExp
            \end{itemize}
        \end{tcolorbox}
    \end{minipage}\hfill
    \begin{minipage}[t]{0.485\textwidth}
        \begin{tcolorbox}[cardAdultoBox, colframe=greenColor!65, coltitle=greenColor!90!black, title={\faTasks\ 6. Tasks (Primary and secondary)}]
            \fontsize{8.0pt}{9.4pt}\selectfont
            \begin{itemize}
                \fAdTasks
            \end{itemize}
        \end{tcolorbox}
    \end{minipage}
    \vspace{0.7mm}
    \noindent\begin{minipage}[t]{0.485\textwidth}
        \begin{tcolorbox}[cardAdultoBox, colframe=accentColor!65, coltitle=accentColor!90!black, title={\faBrain\ 7. Domain knowledge}]
            \fontsize{8.0pt}{9.4pt}\selectfont
            \begin{itemize}
                \fAdDomainKnowledge
            \end{itemize}
        \end{tcolorbox}
    \end{minipage}\hfill
    \begin{minipage}[t]{0.485\textwidth}
        \begin{tcolorbox}[cardAdultoBox, colframe=secondaryColor!60, coltitle=secondaryColor!90!black, title={\faTools\ 8. Technology available}]
            \fontsize{8.0pt}{9.4pt}\selectfont
            \begin{itemize}
                \fAdTechnology
            \end{itemize}
        \end{tcolorbox}
    \end{minipage}
    \vspace{0.7mm}
    \noindent\begin{minipage}[t]{0.485\textwidth}
        \begin{tcolorbox}[cardAdultoBox, colframe=purpleColor!60, coltitle=purpleColor!90!black, title={\faHeart\ 9. Attitudes and values}]
            \fontsize{8.0pt}{9.4pt}\selectfont
            \begin{itemize}
                \fAdAttitudes
            \end{itemize}
        \end{tcolorbox}
    \end{minipage}\hfill
    \begin{minipage}[t]{0.485\textwidth}
        \begin{tcolorbox}[cardAdultoBox, colframe=purpleColor!60, coltitle=purpleColor!90!black, title={\faEye\ 10. Learning style}]
            \fontsize{8.0pt}{9.4pt}\selectfont
            \begin{itemize}
                \fAdLearningStyle
            \end{itemize}
        \end{tcolorbox}
    \end{minipage}
    \vspace{0.7mm}
    % FILA 6: Punto 11 (ANCHO COMPLETO — CRITICALITY OF ERRORS)
    \noindent\begin{minipage}[t]{\textwidth}
        \begin{tcolorbox}[cardAdultoWideBox, colframe=dangerColor!65, coltitle=dangerColor!90!black, title={\faExclamationTriangle\ 11. Criticality of errors (Consecuencias de un error y tolerancia del software)}]
            \fontsize{8.0pt}{9.4pt}\selectfont
            \begin{itemize}
                \fAdCriticality
            \end{itemize}
        \end{tcolorbox}
    \end{minipage}
    \vspace{0.7mm}
    % 4. BARRA INFERIOR DE INTEGRACIÓN CeRA
    \begin{tcolorbox}[
        enhanced,
        colback=tealColor!4,
        colframe=tealColor!45,
        arc=1.6mm,
        left=6pt, right=6pt, top=2.5pt, bottom=2.5pt,
        boxsep=0pt
    ]
        \scriptsize
        \textbf{\color{tealColor!90!black}\faMapSigns\ Escenario de Uso Contextual (Día típico):} \fAdEscenario\\
        \vspace{0.3mm}
        \hrule height 0.35pt \relax
        \vspace{0.5mm}
        \centerline{%
            \textbf{\color{tealColor!90!black}\faSlidersH\ Rol y Funciones en Plataforma (CeRA):} \hspace{2mm}
            \fAdBadgesAjustes%
        }
    \end{tcolorbox}
    \newpage
    \endgroup
}

% COMANDO PARA PLANTILLA EN BLANCO DE ADULTO
\newcommand{\fichaAdultoBlanco}{%
    \fichaAdulto[%
        bannerTitulo={\faUserEdit\ FICHA DE PERSONA: TUTOR/FAMILIA \textbar\ PLANTILLA DE TRABAJO},
        cita={[Cita textual representativa que exprese sus objetivos, inquietudes de gestión o necesidades respecto a la app CeRA].}%
    ]%
}

\ifdefined\isSubFile
    \expandafter\endinput
\fi

\begin{document}
\fichaAdultoBlanco
\end{document}
 y luego invocar \fichaAdulto[...] o \fichaAdultoBlanco
% ==============================================================================

\ifdefined\isSubFile
    % Si se incluye desde otro documento, los paquetes ya están cargados
\else
    \documentclass[10pt]{article}
    \usepackage[utf8]{inputenc}
    \usepackage[spanish,es-tabla]{babel}
    \usepackage[margin=0.8cm,top=0.6cm,bottom=0.6cm]{geometry}
    \usepackage{helvet}
    \renewcommand{\familydefault}{\sfdefault}
    \usepackage{microtype}
    \usepackage{graphicx}
    \usepackage{xcolor}
    \usepackage{tikz}
    \usepackage{enumitem}
    \usepackage{fontawesome5}
    \usepackage{tcolorbox}
    \tcbuselibrary{skins}
    \pagestyle{empty}
\fi

% Colores del sistema
\providecolor{primaryColor}{RGB}{20, 65, 105}     % Azul corporativo
\providecolor{secondaryColor}{RGB}{35, 115, 170}  % Azul medio accesible
\providecolor{tealColor}{RGB}{20, 130, 140}       % Verde azulado
\providecolor{greenColor}{RGB}{30, 135, 75}       % Verde bosque
\providecolor{accentColor}{RGB}{215, 105, 25}     % Naranja cálido
\providecolor{purpleColor}{RGB}{120, 60, 150}     % Púrpura suave
\providecolor{dangerColor}{RGB}{185, 45, 35}      % Rojo criticidad
\providecolor{bgCard}{RGB}{252, 253, 255}         % Fondo limpio
\providecolor{borderCard}{RGB}{215, 225, 235}     % Borde sutil
\providecolor{textDark}{RGB}{30, 40, 50}          % Texto principal oscuro
\providecolor{textMuted}{RGB}{95, 110, 125}       % Texto secundario

% Configuración de listas
\setlist[itemize]{leftmargin=3.0mm, itemsep=0.4pt, parsep=0pt, topsep=0.4pt, partopsep=0pt}

% Estilos de caja fija para Adulto
\tcbset{
    cardAdultoBox/.style={
        enhanced,
        boxrule=0.65pt,
        arc=1.8mm,
        height=27.2mm,
        left=5.0pt, right=5.0pt, top=1.5pt, bottom=1.5pt,
        boxsep=0pt,
        colback=bgCard,
        colframe=borderCard,
        fonttitle=\bfseries\small,
        fontupper=\fontsize{8.2pt}{9.6pt}\selectfont,
    },
    cardAdultoWideBox/.style={
        enhanced,
        boxrule=0.65pt,
        arc=1.8mm,
        height=20.0mm,
        left=6.0pt, right=6.0pt, top=2.0pt, bottom=2.0pt,
        boxsep=0pt,
        colback=bgCard,
        colframe=borderCard,
        fonttitle=\bfseries\small,
    }
}

% Macro para badges visuales (si no existe)
\providecommand{\badge}[2]{%
    \tikz[baseline=(b.base)] \node[fill=#1!12, text=#1!90!black, draw=#1!35, line width=0.4pt, rounded corners=1.2mm, inner sep=2.0pt, font=\bfseries\scriptsize] (b) {#2};%
}

% Definición de claves con pgfkeys para \fichaAdulto
\pgfkeys{
    /fichaAdulto/.is family, /fichaAdulto,
    default/.style = {
        bannerTitulo = {\faChalkboardTeacher\ FICHA DE PERSONA: TUTORES Y FAMILIAS (METODOLOGÍA DCU)},
        subtitulo = {Proyecto AUTOVIDA (26-175) \textbar\ Reto CeRA — Agenda Accesible PTVAL \textbar\ Asignatura MDA (UGR)},
        equipo = {\textbf{Equipo CeRA-MDA} (5 integrantes)\\\ Course 2026/2027 \textbar\ \textbf{Práctica 1}},
        nombre = {[Nombre y Apellidos del Usuario]},
        nombreCorto = {[ADULTO]},
        edad = {[Edad]},
        icono = {\faChalkboardTeacher},
        foto = {},
        badgeRol = {[Rol: Tutor / Familiar / Admin]},
        badgeTarea = {[Ámbito de Actuación]},
        centro = {CEE Purísima Concepción},
        cita = {[Cita textual representativa que exprese sus objetivos, inquietudes de gestión o necesidades respecto a la app CeRA].},
        demographic = {
            \item \textbf{Edad y Género:} [[Edad] años, [Género]].
            \item \textbf{Ubicación:} [[Ciudad / Zona residencial]].
        },
        occupation = {
            \item \textbf{Puesto actual:} [[Puesto docente, profesión o rol familiar]].
            \item \textbf{Antigüedad:} [[Años en el centro/entidad] \textbar\ [Años en puesto]].
            \item \textbf{Responsabilidades:} [[Supervisión, docencia, apoyo en el hogar]].
        },
        education = {
            \item \textbf{Titulación (Degree):} [[Nivel formativo o titulación alcanzada]].
            \item \textbf{Especialidad (Major):} [[Educación Especial, Pedagogía, etc.]].
            \item \textbf{Cursos realizados:} [[Formación en NEAE, SAAC, TIC o gestión]].
        },
        computerExp = {
            \item \textbf{Habilidades informáticas:} [[Nivel de competencia digital: medio/avanzado]].
            \item \textbf{Años de experiencia:} [[Años de uso regular de dispositivos]].
            \item \textbf{Hábitos de uso:} [[Navegación web, apps de gestión, smartphones]].
        },
        productExp = {
            \item \textbf{Productos similares:} [[Otras herramientas usadas, papel, agendas]].
            \item \textbf{Limitaciones actuales:} [[Sobrecarga de trabajo, falta de personalización]].
            \item \textbf{Expectativas:} [[Ahorro de tiempo y canal de dudas]].
        },
        tasks = {
            \item \textbf{Tareas Primarias:} 1) [[Configurar tareas]] $\rightarrow$ 2) [[Asignar a alumnos]] $\rightarrow$ 3) [[Supervisar ejecución]] $\rightarrow$ 4) [[Gestionar comedor/avisos]].
            \item \textbf{Tareas Secundarias:} [[Revisar estadísticas / chat con familias]].
        },
        domainKnowledge = {
            \item \textbf{Comprensión del área:} [[Conocimiento sobre necesidades NEAE y rutinas]].
            \item \textbf{Conocimiento del usuario:} [[Grado de personalización que requiere cada estudiante]].
        },
        technology = {
            \item \textbf{Hardware:} [[Ordenador portátil / tablet / móvil Android]].
            \item \textbf{Software:} [[Navegador web / app móvil]].
            \item \textbf{Herramientas:} [[Conexión a internet]].
        },
        attitudes = {
            \item \textbf{Preferencias:} [[Sencillez, rapidez operativa y fiabilidad en tiempo real]].
            \item \textbf{Miedo a la tecnología:} [[Preocupación por pérdida de datos o fallos de red]].
            \item \textbf{Valores:} [[Foco en el progreso y bienestar del alumnado]].
        },
        learningStyle = {
            \item \textbf{Perfil de aprendizaje:} [[\textbf{Práctico} / \textbf{Estructurado} / \textbf{Visual}]].
            \item \textbf{Canales óptimos:} [[Tutoriales cortos, interfaz intuitiva sin manual]].
            \item \textbf{Curva de adopción:} [[Inmediata si soluciona problemas reales]].
        },
        criticality = {
            \item \textbf{Consecuencias de un error:} \textit{Logístico:} [[¿Dietas erróneas, tareas desincronizadas?]] \textit{Pedagógico:} [[¿Confusión en el horario del alumno, estrés?]]
            \item \textbf{Tolerancia del software:} [[Alertas claras ante cambios sensibles, confirmación de guardado]].
        },
        escenario = {[[Describir brevemente en 3-4 líneas cómo accede al panel de la app, planifica o revisa las tareas de sus estudiantes y cómo la herramienta le facilita la jornada eliminando fricciones organizativas.]]},
        badgesAjustes = {
            \badge{tealColor}{\faDesktop\ [[Panel Web]]} \hspace{2mm}
            \badge{primaryColor}{\faCalendarCheck\ [[Planificación]]} \hspace{2mm}
            \badge{accentColor}{\faUtensils\ [[Comedor/Dietas]]} \hspace{2mm}
            \badge{greenColor}{\faCheckCircle\ [[Supervisión]]}%
        }
    },
    bannerTitulo/.store in = \fAdBannerTitulo,
    subtitulo/.store in = \fAdSubtitulo,
    equipo/.store in = \fAdEquipo,
    nombre/.store in = \fAdNombre,
    nombreCorto/.store in = \fAdNombreCorto,
    edad/.store in = \fAdEdad,
    icono/.store in = \fAdIcono,
    foto/.store in = \fAdFoto,
    badgeRol/.store in = \fAdBadgeRol,
    badgeTarea/.store in = \fAdBadgeTarea,
    centro/.store in = \fAdCentro,
    cita/.store in = \fAdCita,
    demographic/.store in = \fAdDemographic,
    occupation/.store in = \fAdOccupation,
    education/.store in = \fAdEducation,
    computerExp/.store in = \fAdComputerExp,
    productExp/.store in = \fAdProductExp,
    tasks/.store in = \fAdTasks,
    domainKnowledge/.store in = \fAdDomainKnowledge,
    technology/.store in = \fAdTechnology,
    attitudes/.store in = \fAdAttitudes,
    learningStyle/.store in = \fAdLearningStyle,
    criticality/.store in = \fAdCriticality,
    escenario/.store in = \fAdEscenario,
    badgesAjustes/.store in = \fAdBadgesAjustes,
}

% COMANDO PRINCIPAL \fichaAdulto
\newcommand{\fichaAdulto}[1][]{%
    \begingroup
    \pgfkeys{/fichaAdulto, default, #1}%
    % 1. BANNER SUPERIOR
    \begin{tcolorbox}[
        enhanced,
        colback=tealColor!90!black,
        colframe=tealColor!90!black,
        arc=1.6mm,
        left=7pt, right=7pt, top=2.5pt, bottom=2.5pt,
        boxsep=0pt, coltext=white
    ]
        \begin{minipage}[c]{0.70\textwidth}
            {\normalsize\textbf{\fAdBannerTitulo}}\\[0.2mm]
            {\tiny \fAdSubtitulo}
        \end{minipage}\hfill
        \begin{minipage}[c]{0.29\textwidth}
            \raggedleft\tiny
            \fAdEquipo
        \end{minipage}
    \end{tcolorbox}
    \vspace{0.7mm}
    % 2. TARJETA DE IDENTIDAD DEL ADULTO
    \begin{tcolorbox}[
        enhanced,
        colback=white,
        colframe=tealColor!70,
        arc=2.0mm,
        boxrule=0.8pt,
        left=6pt, right=6pt, top=2.8pt, bottom=2.8pt,
        boxsep=0pt,
        drop shadow=black!5
    ]
        \begin{minipage}[c]{0.10\textwidth}
            \centering
            \begin{tikzpicture}
                \ifx\fAdFoto\empty
                    \draw[fill=tealColor!8, draw=tealColor!30, line width=0.5pt, rounded corners=2mm] (0,0) rectangle (1.4, 1.55);
                    \node at (0.7, 0.95) {\color{tealColor}\scalebox{1.5}{\fAdIcono}};
                    \node[align=center] at (0.7, 0.28) {\tiny\textbf{\color{tealColor}\fAdNombreCorto}\\\fontsize{4pt}{5pt}\selectfont\color{textMuted}\fAdEdad};
                \else
                    \begin{scope}
                        \clip[rounded corners=2mm] (0,0) rectangle (1.45, 1.55);
                        \node[inner sep=0pt] at (0.725, 0.775) {\includegraphics[width=1.45cm,height=1.55cm,keepaspectratio=false]{\fAdFoto}};
                    \end{scope}
                    \draw[draw=tealColor!60, line width=0.8pt, rounded corners=2mm] (0,0) rectangle (1.45, 1.55);
                \fi
            \end{tikzpicture}
        \end{minipage}\hfill
        \begin{minipage}[c]{0.88\textwidth}
            \begin{minipage}[b]{0.68\textwidth}\raggedright
                {\large\textbf{\color{tealColor!90!black}\fAdNombre}} \hspace{2mm}
                \badge{tealColor}{\fAdBadgeRol} \badge{secondaryColor}{\fAdBadgeTarea}
            \end{minipage}\hfill
            \begin{minipage}[b]{0.31\textwidth}
                \raggedleft\small\textbf{\color{tealColor!90!black}\faSchool\ \fAdCentro}
            \end{minipage}\\[0.6mm]
            \tikz[baseline] \node[fill=tealColor!5, draw=tealColor!25, rounded corners=1.2mm, inner sep=2.5pt, text width=0.98\linewidth] {
                \small\itshape\color{tealColor!90!black}
                \faQuoteLeft\ ``\fAdCita''\ \faQuoteRight
            };
        \end{minipage}
    \end{tcolorbox}
    \vspace{0.7mm}
    % 3. LAS 11 CAJAS DE ADULTOS (5 FILAS X 2 COLS + 1 FILA COMPLETA)
    \noindent\begin{minipage}[t]{0.485\textwidth}
        \begin{tcolorbox}[cardAdultoBox, colframe=primaryColor!50, coltitle=primaryColor, title={\faIdCard\ 1. Demographic characteristics}]
            \fontsize{8.0pt}{9.4pt}\selectfont
            \begin{itemize}
                \fAdDemographic
            \end{itemize}
        \end{tcolorbox}
    \end{minipage}\hfill
    \begin{minipage}[t]{0.485\textwidth}
        \begin{tcolorbox}[cardAdultoBox, colframe=primaryColor!50, coltitle=primaryColor, title={\faUserTie\ 2. Occupation experience}]
            \fontsize{8.0pt}{9.4pt}\selectfont
            \begin{itemize}
                \fAdOccupation
            \end{itemize}
        \end{tcolorbox}
    \end{minipage}
    \vspace{0.7mm}
    \noindent\begin{minipage}[t]{0.485\textwidth}
        \begin{tcolorbox}[cardAdultoBox, colframe=tealColor!60, coltitle=tealColor!90!black, title={\faGraduationCap\ 3. Education}]
            \fontsize{8.0pt}{9.4pt}\selectfont
            \begin{itemize}
                \fAdEducation
            \end{itemize}
        \end{tcolorbox}
    \end{minipage}\hfill
    \begin{minipage}[t]{0.485\textwidth}
        \begin{tcolorbox}[cardAdultoBox, colframe=secondaryColor!60, coltitle=secondaryColor!90!black, title={\faLaptop\ 4. Computer experience}]
            \fontsize{8.0pt}{9.4pt}\selectfont
            \begin{itemize}
                \fAdComputerExp
            \end{itemize}
        \end{tcolorbox}
    \end{minipage}
    \vspace{0.7mm}
    \noindent\begin{minipage}[t]{0.485\textwidth}
        \begin{tcolorbox}[cardAdultoBox, colframe=secondaryColor!60, coltitle=secondaryColor!90!black, title={\faSearchPlus\ 5. Specific product experience}]
            \fontsize{8.0pt}{9.4pt}\selectfont
            \begin{itemize}
                \fAdProductExp
            \end{itemize}
        \end{tcolorbox}
    \end{minipage}\hfill
    \begin{minipage}[t]{0.485\textwidth}
        \begin{tcolorbox}[cardAdultoBox, colframe=greenColor!65, coltitle=greenColor!90!black, title={\faTasks\ 6. Tasks (Primary and secondary)}]
            \fontsize{8.0pt}{9.4pt}\selectfont
            \begin{itemize}
                \fAdTasks
            \end{itemize}
        \end{tcolorbox}
    \end{minipage}
    \vspace{0.7mm}
    \noindent\begin{minipage}[t]{0.485\textwidth}
        \begin{tcolorbox}[cardAdultoBox, colframe=accentColor!65, coltitle=accentColor!90!black, title={\faBrain\ 7. Domain knowledge}]
            \fontsize{8.0pt}{9.4pt}\selectfont
            \begin{itemize}
                \fAdDomainKnowledge
            \end{itemize}
        \end{tcolorbox}
    \end{minipage}\hfill
    \begin{minipage}[t]{0.485\textwidth}
        \begin{tcolorbox}[cardAdultoBox, colframe=secondaryColor!60, coltitle=secondaryColor!90!black, title={\faTools\ 8. Technology available}]
            \fontsize{8.0pt}{9.4pt}\selectfont
            \begin{itemize}
                \fAdTechnology
            \end{itemize}
        \end{tcolorbox}
    \end{minipage}
    \vspace{0.7mm}
    \noindent\begin{minipage}[t]{0.485\textwidth}
        \begin{tcolorbox}[cardAdultoBox, colframe=purpleColor!60, coltitle=purpleColor!90!black, title={\faHeart\ 9. Attitudes and values}]
            \fontsize{8.0pt}{9.4pt}\selectfont
            \begin{itemize}
                \fAdAttitudes
            \end{itemize}
        \end{tcolorbox}
    \end{minipage}\hfill
    \begin{minipage}[t]{0.485\textwidth}
        \begin{tcolorbox}[cardAdultoBox, colframe=purpleColor!60, coltitle=purpleColor!90!black, title={\faEye\ 10. Learning style}]
            \fontsize{8.0pt}{9.4pt}\selectfont
            \begin{itemize}
                \fAdLearningStyle
            \end{itemize}
        \end{tcolorbox}
    \end{minipage}
    \vspace{0.7mm}
    % FILA 6: Punto 11 (ANCHO COMPLETO — CRITICALITY OF ERRORS)
    \noindent\begin{minipage}[t]{\textwidth}
        \begin{tcolorbox}[cardAdultoWideBox, colframe=dangerColor!65, coltitle=dangerColor!90!black, title={\faExclamationTriangle\ 11. Criticality of errors (Consecuencias de un error y tolerancia del software)}]
            \fontsize{8.0pt}{9.4pt}\selectfont
            \begin{itemize}
                \fAdCriticality
            \end{itemize}
        \end{tcolorbox}
    \end{minipage}
    \vspace{0.7mm}
    % 4. BARRA INFERIOR DE INTEGRACIÓN CeRA
    \begin{tcolorbox}[
        enhanced,
        colback=tealColor!4,
        colframe=tealColor!45,
        arc=1.6mm,
        left=6pt, right=6pt, top=2.5pt, bottom=2.5pt,
        boxsep=0pt
    ]
        \scriptsize
        \textbf{\color{tealColor!90!black}\faMapSigns\ Escenario de Uso Contextual (Día típico):} \fAdEscenario\\
        \vspace{0.3mm}
        \hrule height 0.35pt \relax
        \vspace{0.5mm}
        \centerline{%
            \textbf{\color{tealColor!90!black}\faSlidersH\ Rol y Funciones en Plataforma (CeRA):} \hspace{2mm}
            \fAdBadgesAjustes%
        }
    \end{tcolorbox}
    \newpage
    \endgroup
}

% COMANDO PARA PLANTILLA EN BLANCO DE ADULTO
\newcommand{\fichaAdultoBlanco}{%
    \fichaAdulto[%
        bannerTitulo={\faUserEdit\ FICHA DE PERSONA: TUTOR/FAMILIA \textbar\ PLANTILLA DE TRABAJO},
        cita={[Cita textual representativa que exprese sus objetivos, inquietudes de gestión o necesidades respecto a la app CeRA].}%
    ]%
}

\ifdefined\isSubFile
    \expandafter\endinput
\fi

\begin{document}
\fichaAdultoBlanco
\end{document}
 y luego invocar \fichaAdulto[...] o \fichaAdultoBlanco
% ==============================================================================

\ifdefined\isSubFile
    % Si se incluye desde otro documento, los paquetes ya están cargados
\else
    \documentclass[10pt]{article}
    \usepackage[utf8]{inputenc}
    \usepackage[spanish,es-tabla]{babel}
    \usepackage[margin=0.8cm,top=0.6cm,bottom=0.6cm]{geometry}
    \usepackage{helvet}
    \renewcommand{\familydefault}{\sfdefault}
    \usepackage{microtype}
    \usepackage{graphicx}
    \usepackage{xcolor}
    \usepackage{tikz}
    \usepackage{enumitem}
    \usepackage{fontawesome5}
    \usepackage{tcolorbox}
    \tcbuselibrary{skins}
    \pagestyle{empty}
\fi

% Colores del sistema
\providecolor{primaryColor}{RGB}{20, 65, 105}     % Azul corporativo
\providecolor{secondaryColor}{RGB}{35, 115, 170}  % Azul medio accesible
\providecolor{tealColor}{RGB}{20, 130, 140}       % Verde azulado
\providecolor{greenColor}{RGB}{30, 135, 75}       % Verde bosque
\providecolor{accentColor}{RGB}{215, 105, 25}     % Naranja cálido
\providecolor{purpleColor}{RGB}{120, 60, 150}     % Púrpura suave
\providecolor{dangerColor}{RGB}{185, 45, 35}      % Rojo criticidad
\providecolor{bgCard}{RGB}{252, 253, 255}         % Fondo limpio
\providecolor{borderCard}{RGB}{215, 225, 235}     % Borde sutil
\providecolor{textDark}{RGB}{30, 40, 50}          % Texto principal oscuro
\providecolor{textMuted}{RGB}{95, 110, 125}       % Texto secundario

% Configuración de listas
\setlist[itemize]{leftmargin=3.0mm, itemsep=0.4pt, parsep=0pt, topsep=0.4pt, partopsep=0pt}

% Estilos de caja fija para Adulto
\tcbset{
    cardAdultoBox/.style={
        enhanced,
        boxrule=0.65pt,
        arc=1.8mm,
        height=27.2mm,
        left=5.0pt, right=5.0pt, top=1.5pt, bottom=1.5pt,
        boxsep=0pt,
        colback=bgCard,
        colframe=borderCard,
        fonttitle=\bfseries\small,
        fontupper=\fontsize{8.2pt}{9.6pt}\selectfont,
    },
    cardAdultoWideBox/.style={
        enhanced,
        boxrule=0.65pt,
        arc=1.8mm,
        height=20.0mm,
        left=6.0pt, right=6.0pt, top=2.0pt, bottom=2.0pt,
        boxsep=0pt,
        colback=bgCard,
        colframe=borderCard,
        fonttitle=\bfseries\small,
    }
}

% Macro para badges visuales (si no existe)
\providecommand{\badge}[2]{%
    \tikz[baseline=(b.base)] \node[fill=#1!12, text=#1!90!black, draw=#1!35, line width=0.4pt, rounded corners=1.2mm, inner sep=2.0pt, font=\bfseries\scriptsize] (b) {#2};%
}

% Definición de claves con pgfkeys para \fichaAdulto
\pgfkeys{
    /fichaAdulto/.is family, /fichaAdulto,
    default/.style = {
        bannerTitulo = {\faChalkboardTeacher\ FICHA DE PERSONA: TUTORES Y FAMILIAS (METODOLOGÍA DCU)},
        subtitulo = {Proyecto AUTOVIDA (26-175) \textbar\ Reto CeRA — Agenda Accesible PTVAL \textbar\ Asignatura MDA (UGR)},
        equipo = {\textbf{Equipo CeRA-MDA} (5 integrantes)\\\ Course 2026/2027 \textbar\ \textbf{Práctica 1}},
        nombre = {[Nombre y Apellidos del Usuario]},
        nombreCorto = {[ADULTO]},
        edad = {[Edad]},
        icono = {\faChalkboardTeacher},
        foto = {},
        badgeRol = {[Rol: Tutor / Familiar / Admin]},
        badgeTarea = {[Ámbito de Actuación]},
        centro = {CEE Purísima Concepción},
        cita = {[Cita textual representativa que exprese sus objetivos, inquietudes de gestión o necesidades respecto a la app CeRA].},
        demographic = {
            \item \textbf{Edad y Género:} [[Edad] años, [Género]].
            \item \textbf{Ubicación:} [[Ciudad / Zona residencial]].
        },
        occupation = {
            \item \textbf{Puesto actual:} [[Puesto docente, profesión o rol familiar]].
            \item \textbf{Antigüedad:} [[Años en el centro/entidad] \textbar\ [Años en puesto]].
            \item \textbf{Responsabilidades:} [[Supervisión, docencia, apoyo en el hogar]].
        },
        education = {
            \item \textbf{Titulación (Degree):} [[Nivel formativo o titulación alcanzada]].
            \item \textbf{Especialidad (Major):} [[Educación Especial, Pedagogía, etc.]].
            \item \textbf{Cursos realizados:} [[Formación en NEAE, SAAC, TIC o gestión]].
        },
        computerExp = {
            \item \textbf{Habilidades informáticas:} [[Nivel de competencia digital: medio/avanzado]].
            \item \textbf{Años de experiencia:} [[Años de uso regular de dispositivos]].
            \item \textbf{Hábitos de uso:} [[Navegación web, apps de gestión, smartphones]].
        },
        productExp = {
            \item \textbf{Productos similares:} [[Otras herramientas usadas, papel, agendas]].
            \item \textbf{Limitaciones actuales:} [[Sobrecarga de trabajo, falta de personalización]].
            \item \textbf{Expectativas:} [[Ahorro de tiempo y canal de dudas]].
        },
        tasks = {
            \item \textbf{Tareas Primarias:} 1) [[Configurar tareas]] $\rightarrow$ 2) [[Asignar a alumnos]] $\rightarrow$ 3) [[Supervisar ejecución]] $\rightarrow$ 4) [[Gestionar comedor/avisos]].
            \item \textbf{Tareas Secundarias:} [[Revisar estadísticas / chat con familias]].
        },
        domainKnowledge = {
            \item \textbf{Comprensión del área:} [[Conocimiento sobre necesidades NEAE y rutinas]].
            \item \textbf{Conocimiento del usuario:} [[Grado de personalización que requiere cada estudiante]].
        },
        technology = {
            \item \textbf{Hardware:} [[Ordenador portátil / tablet / móvil Android]].
            \item \textbf{Software:} [[Navegador web / app móvil]].
            \item \textbf{Herramientas:} [[Conexión a internet]].
        },
        attitudes = {
            \item \textbf{Preferencias:} [[Sencillez, rapidez operativa y fiabilidad en tiempo real]].
            \item \textbf{Miedo a la tecnología:} [[Preocupación por pérdida de datos o fallos de red]].
            \item \textbf{Valores:} [[Foco en el progreso y bienestar del alumnado]].
        },
        learningStyle = {
            \item \textbf{Perfil de aprendizaje:} [[\textbf{Práctico} / \textbf{Estructurado} / \textbf{Visual}]].
            \item \textbf{Canales óptimos:} [[Tutoriales cortos, interfaz intuitiva sin manual]].
            \item \textbf{Curva de adopción:} [[Inmediata si soluciona problemas reales]].
        },
        criticality = {
            \item \textbf{Consecuencias de un error:} \textit{Logístico:} [[¿Dietas erróneas, tareas desincronizadas?]] \textit{Pedagógico:} [[¿Confusión en el horario del alumno, estrés?]]
            \item \textbf{Tolerancia del software:} [[Alertas claras ante cambios sensibles, confirmación de guardado]].
        },
        escenario = {[[Describir brevemente en 3-4 líneas cómo accede al panel de la app, planifica o revisa las tareas de sus estudiantes y cómo la herramienta le facilita la jornada eliminando fricciones organizativas.]]},
        badgesAjustes = {
            \badge{tealColor}{\faDesktop\ [[Panel Web]]} \hspace{2mm}
            \badge{primaryColor}{\faCalendarCheck\ [[Planificación]]} \hspace{2mm}
            \badge{accentColor}{\faUtensils\ [[Comedor/Dietas]]} \hspace{2mm}
            \badge{greenColor}{\faCheckCircle\ [[Supervisión]]}%
        }
    },
    bannerTitulo/.store in = \fAdBannerTitulo,
    subtitulo/.store in = \fAdSubtitulo,
    equipo/.store in = \fAdEquipo,
    nombre/.store in = \fAdNombre,
    nombreCorto/.store in = \fAdNombreCorto,
    edad/.store in = \fAdEdad,
    icono/.store in = \fAdIcono,
    foto/.store in = \fAdFoto,
    badgeRol/.store in = \fAdBadgeRol,
    badgeTarea/.store in = \fAdBadgeTarea,
    centro/.store in = \fAdCentro,
    cita/.store in = \fAdCita,
    demographic/.store in = \fAdDemographic,
    occupation/.store in = \fAdOccupation,
    education/.store in = \fAdEducation,
    computerExp/.store in = \fAdComputerExp,
    productExp/.store in = \fAdProductExp,
    tasks/.store in = \fAdTasks,
    domainKnowledge/.store in = \fAdDomainKnowledge,
    technology/.store in = \fAdTechnology,
    attitudes/.store in = \fAdAttitudes,
    learningStyle/.store in = \fAdLearningStyle,
    criticality/.store in = \fAdCriticality,
    escenario/.store in = \fAdEscenario,
    badgesAjustes/.store in = \fAdBadgesAjustes,
}

% COMANDO PRINCIPAL \fichaAdulto
\newcommand{\fichaAdulto}[1][]{%
    \begingroup
    \pgfkeys{/fichaAdulto, default, #1}%
    % 1. BANNER SUPERIOR
    \begin{tcolorbox}[
        enhanced,
        colback=tealColor!90!black,
        colframe=tealColor!90!black,
        arc=1.6mm,
        left=7pt, right=7pt, top=2.5pt, bottom=2.5pt,
        boxsep=0pt, coltext=white
    ]
        \begin{minipage}[c]{0.70\textwidth}
            {\normalsize\textbf{\fAdBannerTitulo}}\\[0.2mm]
            {\tiny \fAdSubtitulo}
        \end{minipage}\hfill
        \begin{minipage}[c]{0.29\textwidth}
            \raggedleft\tiny
            \fAdEquipo
        \end{minipage}
    \end{tcolorbox}
    \vspace{0.7mm}
    % 2. TARJETA DE IDENTIDAD DEL ADULTO
    \begin{tcolorbox}[
        enhanced,
        colback=white,
        colframe=tealColor!70,
        arc=2.0mm,
        boxrule=0.8pt,
        left=6pt, right=6pt, top=2.8pt, bottom=2.8pt,
        boxsep=0pt,
        drop shadow=black!5
    ]
        \begin{minipage}[c]{0.10\textwidth}
            \centering
            \begin{tikzpicture}
                \ifx\fAdFoto\empty
                    \draw[fill=tealColor!8, draw=tealColor!30, line width=0.5pt, rounded corners=2mm] (0,0) rectangle (1.4, 1.55);
                    \node at (0.7, 0.95) {\color{tealColor}\scalebox{1.5}{\fAdIcono}};
                    \node[align=center] at (0.7, 0.28) {\tiny\textbf{\color{tealColor}\fAdNombreCorto}\\\fontsize{4pt}{5pt}\selectfont\color{textMuted}\fAdEdad};
                \else
                    \begin{scope}
                        \clip[rounded corners=2mm] (0,0) rectangle (1.45, 1.55);
                        \node[inner sep=0pt] at (0.725, 0.775) {\includegraphics[width=1.45cm,height=1.55cm,keepaspectratio=false]{\fAdFoto}};
                    \end{scope}
                    \draw[draw=tealColor!60, line width=0.8pt, rounded corners=2mm] (0,0) rectangle (1.45, 1.55);
                \fi
            \end{tikzpicture}
        \end{minipage}\hfill
        \begin{minipage}[c]{0.88\textwidth}
            \begin{minipage}[b]{0.68\textwidth}\raggedright
                {\large\textbf{\color{tealColor!90!black}\fAdNombre}} \hspace{2mm}
                \badge{tealColor}{\fAdBadgeRol} \badge{secondaryColor}{\fAdBadgeTarea}
            \end{minipage}\hfill
            \begin{minipage}[b]{0.31\textwidth}
                \raggedleft\small\textbf{\color{tealColor!90!black}\faSchool\ \fAdCentro}
            \end{minipage}\\[0.6mm]
            \tikz[baseline] \node[fill=tealColor!5, draw=tealColor!25, rounded corners=1.2mm, inner sep=2.5pt, text width=0.98\linewidth] {
                \small\itshape\color{tealColor!90!black}
                \faQuoteLeft\ ``\fAdCita''\ \faQuoteRight
            };
        \end{minipage}
    \end{tcolorbox}
    \vspace{0.7mm}
    % 3. LAS 11 CAJAS DE ADULTOS (5 FILAS X 2 COLS + 1 FILA COMPLETA)
    \noindent\begin{minipage}[t]{0.485\textwidth}
        \begin{tcolorbox}[cardAdultoBox, colframe=primaryColor!50, coltitle=primaryColor, title={\faIdCard\ 1. Demographic characteristics}]
            \fontsize{8.0pt}{9.4pt}\selectfont
            \begin{itemize}
                \fAdDemographic
            \end{itemize}
        \end{tcolorbox}
    \end{minipage}\hfill
    \begin{minipage}[t]{0.485\textwidth}
        \begin{tcolorbox}[cardAdultoBox, colframe=primaryColor!50, coltitle=primaryColor, title={\faUserTie\ 2. Occupation experience}]
            \fontsize{8.0pt}{9.4pt}\selectfont
            \begin{itemize}
                \fAdOccupation
            \end{itemize}
        \end{tcolorbox}
    \end{minipage}
    \vspace{0.7mm}
    \noindent\begin{minipage}[t]{0.485\textwidth}
        \begin{tcolorbox}[cardAdultoBox, colframe=tealColor!60, coltitle=tealColor!90!black, title={\faGraduationCap\ 3. Education}]
            \fontsize{8.0pt}{9.4pt}\selectfont
            \begin{itemize}
                \fAdEducation
            \end{itemize}
        \end{tcolorbox}
    \end{minipage}\hfill
    \begin{minipage}[t]{0.485\textwidth}
        \begin{tcolorbox}[cardAdultoBox, colframe=secondaryColor!60, coltitle=secondaryColor!90!black, title={\faLaptop\ 4. Computer experience}]
            \fontsize{8.0pt}{9.4pt}\selectfont
            \begin{itemize}
                \fAdComputerExp
            \end{itemize}
        \end{tcolorbox}
    \end{minipage}
    \vspace{0.7mm}
    \noindent\begin{minipage}[t]{0.485\textwidth}
        \begin{tcolorbox}[cardAdultoBox, colframe=secondaryColor!60, coltitle=secondaryColor!90!black, title={\faSearchPlus\ 5. Specific product experience}]
            \fontsize{8.0pt}{9.4pt}\selectfont
            \begin{itemize}
                \fAdProductExp
            \end{itemize}
        \end{tcolorbox}
    \end{minipage}\hfill
    \begin{minipage}[t]{0.485\textwidth}
        \begin{tcolorbox}[cardAdultoBox, colframe=greenColor!65, coltitle=greenColor!90!black, title={\faTasks\ 6. Tasks (Primary and secondary)}]
            \fontsize{8.0pt}{9.4pt}\selectfont
            \begin{itemize}
                \fAdTasks
            \end{itemize}
        \end{tcolorbox}
    \end{minipage}
    \vspace{0.7mm}
    \noindent\begin{minipage}[t]{0.485\textwidth}
        \begin{tcolorbox}[cardAdultoBox, colframe=accentColor!65, coltitle=accentColor!90!black, title={\faBrain\ 7. Domain knowledge}]
            \fontsize{8.0pt}{9.4pt}\selectfont
            \begin{itemize}
                \fAdDomainKnowledge
            \end{itemize}
        \end{tcolorbox}
    \end{minipage}\hfill
    \begin{minipage}[t]{0.485\textwidth}
        \begin{tcolorbox}[cardAdultoBox, colframe=secondaryColor!60, coltitle=secondaryColor!90!black, title={\faTools\ 8. Technology available}]
            \fontsize{8.0pt}{9.4pt}\selectfont
            \begin{itemize}
                \fAdTechnology
            \end{itemize}
        \end{tcolorbox}
    \end{minipage}
    \vspace{0.7mm}
    \noindent\begin{minipage}[t]{0.485\textwidth}
        \begin{tcolorbox}[cardAdultoBox, colframe=purpleColor!60, coltitle=purpleColor!90!black, title={\faHeart\ 9. Attitudes and values}]
            \fontsize{8.0pt}{9.4pt}\selectfont
            \begin{itemize}
                \fAdAttitudes
            \end{itemize}
        \end{tcolorbox}
    \end{minipage}\hfill
    \begin{minipage}[t]{0.485\textwidth}
        \begin{tcolorbox}[cardAdultoBox, colframe=purpleColor!60, coltitle=purpleColor!90!black, title={\faEye\ 10. Learning style}]
            \fontsize{8.0pt}{9.4pt}\selectfont
            \begin{itemize}
                \fAdLearningStyle
            \end{itemize}
        \end{tcolorbox}
    \end{minipage}
    \vspace{0.7mm}
    % FILA 6: Punto 11 (ANCHO COMPLETO — CRITICALITY OF ERRORS)
    \noindent\begin{minipage}[t]{\textwidth}
        \begin{tcolorbox}[cardAdultoWideBox, colframe=dangerColor!65, coltitle=dangerColor!90!black, title={\faExclamationTriangle\ 11. Criticality of errors (Consecuencias de un error y tolerancia del software)}]
            \fontsize{8.0pt}{9.4pt}\selectfont
            \begin{itemize}
                \fAdCriticality
            \end{itemize}
        \end{tcolorbox}
    \end{minipage}
    \vspace{0.7mm}
    % 4. BARRA INFERIOR DE INTEGRACIÓN CeRA
    \begin{tcolorbox}[
        enhanced,
        colback=tealColor!4,
        colframe=tealColor!45,
        arc=1.6mm,
        left=6pt, right=6pt, top=2.5pt, bottom=2.5pt,
        boxsep=0pt
    ]
        \scriptsize
        \textbf{\color{tealColor!90!black}\faMapSigns\ Escenario de Uso Contextual (Día típico):} \fAdEscenario\\
        \vspace{0.3mm}
        \hrule height 0.35pt \relax
        \vspace{0.5mm}
        \centerline{%
            \textbf{\color{tealColor!90!black}\faSlidersH\ Rol y Funciones en Plataforma (CeRA):} \hspace{2mm}
            \fAdBadgesAjustes%
        }
    \end{tcolorbox}
    \newpage
    \endgroup
}

% COMANDO PARA PLANTILLA EN BLANCO DE ADULTO
\newcommand{\fichaAdultoBlanco}{%
    \fichaAdulto[%
        bannerTitulo={\faUserEdit\ FICHA DE PERSONA: TUTOR/FAMILIA \textbar\ PLANTILLA DE TRABAJO},
        cita={[Cita textual representativa que exprese sus objetivos, inquietudes de gestión o necesidades respecto a la app CeRA].}%
    ]%
}

\ifdefined\isSubFile
    \expandafter\endinput
\fi

\begin{document}
\fichaAdultoBlanco
\end{document}
 y luego invocar \fichaAdulto[...] o \fichaAdultoBlanco
% ==============================================================================

\ifdefined\isSubFile
    % Si se incluye desde otro documento, los paquetes ya están cargados
\else
    \documentclass[10pt]{article}
    \usepackage[utf8]{inputenc}
    \usepackage[spanish,es-tabla]{babel}
    \usepackage[margin=0.8cm,top=0.6cm,bottom=0.6cm]{geometry}
    \usepackage{helvet}
    \renewcommand{\familydefault}{\sfdefault}
    \usepackage{microtype}
    \usepackage{graphicx}
    \usepackage{xcolor}
    \usepackage{tikz}
    \usepackage{enumitem}
    \usepackage{fontawesome5}
    \usepackage{tcolorbox}
    \tcbuselibrary{skins}
    \pagestyle{empty}
\fi

% Colores del sistema
\providecolor{primaryColor}{RGB}{20, 65, 105}     % Azul corporativo
\providecolor{secondaryColor}{RGB}{35, 115, 170}  % Azul medio accesible
\providecolor{tealColor}{RGB}{20, 130, 140}       % Verde azulado
\providecolor{greenColor}{RGB}{30, 135, 75}       % Verde bosque
\providecolor{accentColor}{RGB}{215, 105, 25}     % Naranja cálido
\providecolor{purpleColor}{RGB}{120, 60, 150}     % Púrpura suave
\providecolor{dangerColor}{RGB}{185, 45, 35}      % Rojo criticidad
\providecolor{bgCard}{RGB}{252, 253, 255}         % Fondo limpio
\providecolor{borderCard}{RGB}{215, 225, 235}     % Borde sutil
\providecolor{textDark}{RGB}{30, 40, 50}          % Texto principal oscuro
\providecolor{textMuted}{RGB}{95, 110, 125}       % Texto secundario

% Configuración de listas
\setlist[itemize]{leftmargin=3.0mm, itemsep=0.4pt, parsep=0pt, topsep=0.4pt, partopsep=0pt}

% Estilos de caja fija para Adulto
\tcbset{
    cardAdultoBox/.style={
        enhanced,
        boxrule=0.65pt,
        arc=1.8mm,
        height=27.2mm,
        left=5.0pt, right=5.0pt, top=1.5pt, bottom=1.5pt,
        boxsep=0pt,
        colback=bgCard,
        colframe=borderCard,
        fonttitle=\bfseries\small,
        fontupper=\fontsize{8.2pt}{9.6pt}\selectfont,
    },
    cardAdultoWideBox/.style={
        enhanced,
        boxrule=0.65pt,
        arc=1.8mm,
        height=20.0mm,
        left=6.0pt, right=6.0pt, top=2.0pt, bottom=2.0pt,
        boxsep=0pt,
        colback=bgCard,
        colframe=borderCard,
        fonttitle=\bfseries\small,
    }
}

% Macro para badges visuales (si no existe)
\providecommand{\badge}[2]{%
    \tikz[baseline=(b.base)] \node[fill=#1!12, text=#1!90!black, draw=#1!35, line width=0.4pt, rounded corners=1.2mm, inner sep=2.0pt, font=\bfseries\scriptsize] (b) {#2};%
}

% Definición de claves con pgfkeys para \fichaAdulto
\pgfkeys{
    /fichaAdulto/.is family, /fichaAdulto,
    default/.style = {
        bannerTitulo = {\faChalkboardTeacher\ FICHA DE PERSONA: TUTORES Y FAMILIAS (METODOLOGÍA DCU)},
        subtitulo = {Proyecto AUTOVIDA (26-175) \textbar\ Reto CeRA — Agenda Accesible PTVAL \textbar\ Asignatura MDA (UGR)},
        equipo = {\textbf{Equipo CeRA-MDA} (5 integrantes)\\\ Course 2026/2027 \textbar\ \textbf{Práctica 1}},
        nombre = {[Nombre y Apellidos del Usuario]},
        nombreCorto = {[ADULTO]},
        edad = {[Edad]},
        icono = {\faChalkboardTeacher},
        foto = {},
        badgeRol = {[Rol: Tutor / Familiar / Admin]},
        badgeTarea = {[Ámbito de Actuación]},
        centro = {CEE Purísima Concepción},
        cita = {[Cita textual representativa que exprese sus objetivos, inquietudes de gestión o necesidades respecto a la app CeRA].},
        demographic = {
            \item \textbf{Edad y Género:} [[Edad] años, [Género]].
            \item \textbf{Ubicación:} [[Ciudad / Zona residencial]].
        },
        occupation = {
            \item \textbf{Puesto actual:} [[Puesto docente, profesión o rol familiar]].
            \item \textbf{Antigüedad:} [[Años en el centro/entidad] \textbar\ [Años en puesto]].
            \item \textbf{Responsabilidades:} [[Supervisión, docencia, apoyo en el hogar]].
        },
        education = {
            \item \textbf{Titulación (Degree):} [[Nivel formativo o titulación alcanzada]].
            \item \textbf{Especialidad (Major):} [[Educación Especial, Pedagogía, etc.]].
            \item \textbf{Cursos realizados:} [[Formación en NEAE, SAAC, TIC o gestión]].
        },
        computerExp = {
            \item \textbf{Habilidades informáticas:} [[Nivel de competencia digital: medio/avanzado]].
            \item \textbf{Años de experiencia:} [[Años de uso regular de dispositivos]].
            \item \textbf{Hábitos de uso:} [[Navegación web, apps de gestión, smartphones]].
        },
        productExp = {
            \item \textbf{Productos similares:} [[Otras herramientas usadas, papel, agendas]].
            \item \textbf{Limitaciones actuales:} [[Sobrecarga de trabajo, falta de personalización]].
            \item \textbf{Expectativas:} [[Ahorro de tiempo y canal de dudas]].
        },
        tasks = {
            \item \textbf{Tareas Primarias:} 1) [[Configurar tareas]] $\rightarrow$ 2) [[Asignar a alumnos]] $\rightarrow$ 3) [[Supervisar ejecución]] $\rightarrow$ 4) [[Gestionar comedor/avisos]].
            \item \textbf{Tareas Secundarias:} [[Revisar estadísticas / chat con familias]].
        },
        domainKnowledge = {
            \item \textbf{Comprensión del área:} [[Conocimiento sobre necesidades NEAE y rutinas]].
            \item \textbf{Conocimiento del usuario:} [[Grado de personalización que requiere cada estudiante]].
        },
        technology = {
            \item \textbf{Hardware:} [[Ordenador portátil / tablet / móvil Android]].
            \item \textbf{Software:} [[Navegador web / app móvil]].
            \item \textbf{Herramientas:} [[Conexión a internet]].
        },
        attitudes = {
            \item \textbf{Preferencias:} [[Sencillez, rapidez operativa y fiabilidad en tiempo real]].
            \item \textbf{Miedo a la tecnología:} [[Preocupación por pérdida de datos o fallos de red]].
            \item \textbf{Valores:} [[Foco en el progreso y bienestar del alumnado]].
        },
        learningStyle = {
            \item \textbf{Perfil de aprendizaje:} [[\textbf{Práctico} / \textbf{Estructurado} / \textbf{Visual}]].
            \item \textbf{Canales óptimos:} [[Tutoriales cortos, interfaz intuitiva sin manual]].
            \item \textbf{Curva de adopción:} [[Inmediata si soluciona problemas reales]].
        },
        criticality = {
            \item \textbf{Consecuencias de un error:} \textit{Logístico:} [[¿Dietas erróneas, tareas desincronizadas?]] \textit{Pedagógico:} [[¿Confusión en el horario del alumno, estrés?]]
            \item \textbf{Tolerancia del software:} [[Alertas claras ante cambios sensibles, confirmación de guardado]].
        },
        escenario = {[[Describir brevemente en 3-4 líneas cómo accede al panel de la app, planifica o revisa las tareas de sus estudiantes y cómo la herramienta le facilita la jornada eliminando fricciones organizativas.]]},
        badgesAjustes = {
            \badge{tealColor}{\faDesktop\ [[Panel Web]]} \hspace{2mm}
            \badge{primaryColor}{\faCalendarCheck\ [[Planificación]]} \hspace{2mm}
            \badge{accentColor}{\faUtensils\ [[Comedor/Dietas]]} \hspace{2mm}
            \badge{greenColor}{\faCheckCircle\ [[Supervisión]]}%
        }
    },
    bannerTitulo/.store in = \fAdBannerTitulo,
    subtitulo/.store in = \fAdSubtitulo,
    equipo/.store in = \fAdEquipo,
    nombre/.store in = \fAdNombre,
    nombreCorto/.store in = \fAdNombreCorto,
    edad/.store in = \fAdEdad,
    icono/.store in = \fAdIcono,
    foto/.store in = \fAdFoto,
    badgeRol/.store in = \fAdBadgeRol,
    badgeTarea/.store in = \fAdBadgeTarea,
    centro/.store in = \fAdCentro,
    cita/.store in = \fAdCita,
    demographic/.store in = \fAdDemographic,
    occupation/.store in = \fAdOccupation,
    education/.store in = \fAdEducation,
    computerExp/.store in = \fAdComputerExp,
    productExp/.store in = \fAdProductExp,
    tasks/.store in = \fAdTasks,
    domainKnowledge/.store in = \fAdDomainKnowledge,
    technology/.store in = \fAdTechnology,
    attitudes/.store in = \fAdAttitudes,
    learningStyle/.store in = \fAdLearningStyle,
    criticality/.store in = \fAdCriticality,
    escenario/.store in = \fAdEscenario,
    badgesAjustes/.store in = \fAdBadgesAjustes,
}

% COMANDO PRINCIPAL \fichaAdulto
\newcommand{\fichaAdulto}[1][]{%
    \begingroup
    \pgfkeys{/fichaAdulto, default, #1}%
    % 1. BANNER SUPERIOR
    \begin{tcolorbox}[
        enhanced,
        colback=tealColor!90!black,
        colframe=tealColor!90!black,
        arc=1.6mm,
        left=7pt, right=7pt, top=2.5pt, bottom=2.5pt,
        boxsep=0pt, coltext=white
    ]
        \begin{minipage}[c]{0.70\textwidth}
            {\normalsize\textbf{\fAdBannerTitulo}}\\[0.2mm]
            {\tiny \fAdSubtitulo}
        \end{minipage}\hfill
        \begin{minipage}[c]{0.29\textwidth}
            \raggedleft\tiny
            \fAdEquipo
        \end{minipage}
    \end{tcolorbox}
    \vspace{0.7mm}
    % 2. TARJETA DE IDENTIDAD DEL ADULTO
    \begin{tcolorbox}[
        enhanced,
        colback=white,
        colframe=tealColor!70,
        arc=2.0mm,
        boxrule=0.8pt,
        left=6pt, right=6pt, top=2.8pt, bottom=2.8pt,
        boxsep=0pt,
        drop shadow=black!5
    ]
        \begin{minipage}[c]{0.10\textwidth}
            \centering
            \begin{tikzpicture}
                \ifx\fAdFoto\empty
                    \draw[fill=tealColor!8, draw=tealColor!30, line width=0.5pt, rounded corners=2mm] (0,0) rectangle (1.4, 1.55);
                    \node at (0.7, 0.95) {\color{tealColor}\scalebox{1.5}{\fAdIcono}};
                    \node[align=center] at (0.7, 0.28) {\tiny\textbf{\color{tealColor}\fAdNombreCorto}\\\fontsize{4pt}{5pt}\selectfont\color{textMuted}\fAdEdad};
                \else
                    \begin{scope}
                        \clip[rounded corners=2mm] (0,0) rectangle (1.45, 1.55);
                        \node[inner sep=0pt] at (0.725, 0.775) {\includegraphics[width=1.45cm,height=1.55cm,keepaspectratio=false]{\fAdFoto}};
                    \end{scope}
                    \draw[draw=tealColor!60, line width=0.8pt, rounded corners=2mm] (0,0) rectangle (1.45, 1.55);
                \fi
            \end{tikzpicture}
        \end{minipage}\hfill
        \begin{minipage}[c]{0.88\textwidth}
            \begin{minipage}[b]{0.68\textwidth}\raggedright
                {\large\textbf{\color{tealColor!90!black}\fAdNombre}} \hspace{2mm}
                \badge{tealColor}{\fAdBadgeRol} \badge{secondaryColor}{\fAdBadgeTarea}
            \end{minipage}\hfill
            \begin{minipage}[b]{0.31\textwidth}
                \raggedleft\small\textbf{\color{tealColor!90!black}\faSchool\ \fAdCentro}
            \end{minipage}\\[0.6mm]
            \tikz[baseline] \node[fill=tealColor!5, draw=tealColor!25, rounded corners=1.2mm, inner sep=2.5pt, text width=0.98\linewidth] {
                \small\itshape\color{tealColor!90!black}
                \faQuoteLeft\ ``\fAdCita''\ \faQuoteRight
            };
        \end{minipage}
    \end{tcolorbox}
    \vspace{0.7mm}
    % 3. LAS 11 CAJAS DE ADULTOS (5 FILAS X 2 COLS + 1 FILA COMPLETA)
    \noindent\begin{minipage}[t]{0.485\textwidth}
        \begin{tcolorbox}[cardAdultoBox, colframe=primaryColor!50, coltitle=primaryColor, title={\faIdCard\ 1. Demographic characteristics}]
            \fontsize{8.0pt}{9.4pt}\selectfont
            \begin{itemize}
                \fAdDemographic
            \end{itemize}
        \end{tcolorbox}
    \end{minipage}\hfill
    \begin{minipage}[t]{0.485\textwidth}
        \begin{tcolorbox}[cardAdultoBox, colframe=primaryColor!50, coltitle=primaryColor, title={\faUserTie\ 2. Occupation experience}]
            \fontsize{8.0pt}{9.4pt}\selectfont
            \begin{itemize}
                \fAdOccupation
            \end{itemize}
        \end{tcolorbox}
    \end{minipage}
    \vspace{0.7mm}
    \noindent\begin{minipage}[t]{0.485\textwidth}
        \begin{tcolorbox}[cardAdultoBox, colframe=tealColor!60, coltitle=tealColor!90!black, title={\faGraduationCap\ 3. Education}]
            \fontsize{8.0pt}{9.4pt}\selectfont
            \begin{itemize}
                \fAdEducation
            \end{itemize}
        \end{tcolorbox}
    \end{minipage}\hfill
    \begin{minipage}[t]{0.485\textwidth}
        \begin{tcolorbox}[cardAdultoBox, colframe=secondaryColor!60, coltitle=secondaryColor!90!black, title={\faLaptop\ 4. Computer experience}]
            \fontsize{8.0pt}{9.4pt}\selectfont
            \begin{itemize}
                \fAdComputerExp
            \end{itemize}
        \end{tcolorbox}
    \end{minipage}
    \vspace{0.7mm}
    \noindent\begin{minipage}[t]{0.485\textwidth}
        \begin{tcolorbox}[cardAdultoBox, colframe=secondaryColor!60, coltitle=secondaryColor!90!black, title={\faSearchPlus\ 5. Specific product experience}]
            \fontsize{8.0pt}{9.4pt}\selectfont
            \begin{itemize}
                \fAdProductExp
            \end{itemize}
        \end{tcolorbox}
    \end{minipage}\hfill
    \begin{minipage}[t]{0.485\textwidth}
        \begin{tcolorbox}[cardAdultoBox, colframe=greenColor!65, coltitle=greenColor!90!black, title={\faTasks\ 6. Tasks (Primary and secondary)}]
            \fontsize{8.0pt}{9.4pt}\selectfont
            \begin{itemize}
                \fAdTasks
            \end{itemize}
        \end{tcolorbox}
    \end{minipage}
    \vspace{0.7mm}
    \noindent\begin{minipage}[t]{0.485\textwidth}
        \begin{tcolorbox}[cardAdultoBox, colframe=accentColor!65, coltitle=accentColor!90!black, title={\faBrain\ 7. Domain knowledge}]
            \fontsize{8.0pt}{9.4pt}\selectfont
            \begin{itemize}
                \fAdDomainKnowledge
            \end{itemize}
        \end{tcolorbox}
    \end{minipage}\hfill
    \begin{minipage}[t]{0.485\textwidth}
        \begin{tcolorbox}[cardAdultoBox, colframe=secondaryColor!60, coltitle=secondaryColor!90!black, title={\faTools\ 8. Technology available}]
            \fontsize{8.0pt}{9.4pt}\selectfont
            \begin{itemize}
                \fAdTechnology
            \end{itemize}
        \end{tcolorbox}
    \end{minipage}
    \vspace{0.7mm}
    \noindent\begin{minipage}[t]{0.485\textwidth}
        \begin{tcolorbox}[cardAdultoBox, colframe=purpleColor!60, coltitle=purpleColor!90!black, title={\faHeart\ 9. Attitudes and values}]
            \fontsize{8.0pt}{9.4pt}\selectfont
            \begin{itemize}
                \fAdAttitudes
            \end{itemize}
        \end{tcolorbox}
    \end{minipage}\hfill
    \begin{minipage}[t]{0.485\textwidth}
        \begin{tcolorbox}[cardAdultoBox, colframe=purpleColor!60, coltitle=purpleColor!90!black, title={\faEye\ 10. Learning style}]
            \fontsize{8.0pt}{9.4pt}\selectfont
            \begin{itemize}
                \fAdLearningStyle
            \end{itemize}
        \end{tcolorbox}
    \end{minipage}
    \vspace{0.7mm}
    % FILA 6: Punto 11 (ANCHO COMPLETO — CRITICALITY OF ERRORS)
    \noindent\begin{minipage}[t]{\textwidth}
        \begin{tcolorbox}[cardAdultoWideBox, colframe=dangerColor!65, coltitle=dangerColor!90!black, title={\faExclamationTriangle\ 11. Criticality of errors (Consecuencias de un error y tolerancia del software)}]
            \fontsize{8.0pt}{9.4pt}\selectfont
            \begin{itemize}
                \fAdCriticality
            \end{itemize}
        \end{tcolorbox}
    \end{minipage}
    \vspace{0.7mm}
    % 4. BARRA INFERIOR DE INTEGRACIÓN CeRA
    \begin{tcolorbox}[
        enhanced,
        colback=tealColor!4,
        colframe=tealColor!45,
        arc=1.6mm,
        left=6pt, right=6pt, top=2.5pt, bottom=2.5pt,
        boxsep=0pt
    ]
        \scriptsize
        \textbf{\color{tealColor!90!black}\faMapSigns\ Escenario de Uso Contextual (Día típico):} \fAdEscenario\\
        \vspace{0.3mm}
        \hrule height 0.35pt \relax
        \vspace{0.5mm}
        \centerline{%
            \textbf{\color{tealColor!90!black}\faSlidersH\ Rol y Funciones en Plataforma (CeRA):} \hspace{2mm}
            \fAdBadgesAjustes%
        }
    \end{tcolorbox}
    \newpage
    \endgroup
}

% COMANDO PARA PLANTILLA EN BLANCO DE ADULTO
\newcommand{\fichaAdultoBlanco}{%
    \fichaAdulto[%
        bannerTitulo={\faUserEdit\ FICHA DE PERSONA: TUTOR/FAMILIA \textbar\ PLANTILLA DE TRABAJO},
        cita={[Cita textual representativa que exprese sus objetivos, inquietudes de gestión o necesidades respecto a la app CeRA].}%
    ]%
}

\ifdefined\isSubFile
    \expandafter\endinput
\fi

\begin{document}
\fichaAdultoBlanco
\end{document}
